\documentclass{amsart}
% For dealing with references we use the comment environment
\usepackage{verbatim}
\newenvironment{reference}{\comment}{\endcomment}
%\newenvironment{reference}{}{}
\newenvironment{slogan}{\comment}{\endcomment}
\newenvironment{history}{\comment}{\endcomment}

% For commutative diagrams we use Xy-pic
\usepackage[all]{xy}

% We use 2cell for 2-commutative diagrams.
\xyoption{2cell}
\UseAllTwocells

% We use multicol for the list of chapters between chapters
\usepackage{multicol}

% This is generally recommended for better output
\usepackage{lmodern}
\usepackage[T1]{fontenc}

% For cross-file-references

% Package for hypertext links:
\usepackage{hyperref}

% For any local file, say "hello.tex" you want to link to please

% Theorem environments.
%
\theoremstyle{plain}
\newtheorem{theorem}[subsection]{Theorem}
\newtheorem{proposition}[subsection]{Proposition}
\newtheorem{lemma}[subsection]{Lemma}

\theoremstyle{definition}
\newtheorem{definition}[subsection]{Definition}
\newtheorem{example}[subsection]{Example}
\newtheorem{exercise}[subsection]{Exercise}
\newtheorem{situation}[subsection]{Situation}

\theoremstyle{remark}
\newtheorem{remark}[subsection]{Remark}
\newtheorem{remarks}[subsection]{Remarks}

\numberwithin{equation}{subsection}

% Macros
%
\def\lim{\mathop{\mathrm{lim}}\nolimits}
\def\colim{\mathop{\mathrm{colim}}\nolimits}
\def\Spec{\mathop{\mathrm{Spec}}}
\def\Hom{\mathop{\mathrm{Hom}}\nolimits}
\def\Ext{\mathop{\mathrm{Ext}}\nolimits}
\def\SheafHom{\mathop{\mathcal{H}\!\mathit{om}}\nolimits}
\def\SheafExt{\mathop{\mathcal{E}\!\mathit{xt}}\nolimits}
\def\Sch{\mathit{Sch}}
\def\Mor{\mathop{\mathrm{Mor}}\nolimits}
\def\Ob{\mathop{\mathrm{Ob}}\nolimits}
\def\Sh{\mathop{\mathit{Sh}}\nolimits}
\def\NL{\mathop{N\!L}\nolimits}
\def\CH{\mathop{\mathrm{CH}}\nolimits}
\def\proetale{{pro\text{-}\acute{e}tale}}
\def\etale{{\acute{e}tale}}
\def\QCoh{\mathit{QCoh}}
\def\Ker{\mathop{\mathrm{Ker}}}
\def\Im{\mathop{\mathrm{Im}}}
\def\Coker{\mathop{\mathrm{Coker}}}
\def\Coim{\mathop{\mathrm{Coim}}}

% Boxtimes
%
\DeclareMathSymbol{\boxtimes}{\mathbin}{AMSa}{"02}

%
% Macros for moduli stacks/spaces
%
\def\QCohstack{\mathcal{QC}\!\mathit{oh}}
\def\Cohstack{\mathcal{C}\!\mathit{oh}}
\def\Spacesstack{\mathcal{S}\!\mathit{paces}}
\def\Quotfunctor{\mathrm{Quot}}
\def\Hilbfunctor{\mathrm{Hilb}}
\def\Curvesstack{\mathcal{C}\!\mathit{urves}}
\def\Polarizedstack{\mathcal{P}\!\mathit{olarized}}
\def\Complexesstack{\mathcal{C}\!\mathit{omplexes}}
% \Pic is the operator that assigns to X its picard group, usage \Pic(X)
% \Picardstack_{X/B} denotes the Picard stack of X over B
% \Picardfunctor_{X/B} denotes the Picard functor of X over B
\def\Pic{\mathop{\mathrm{Pic}}\nolimits}
\def\Picardstack{\mathcal{P}\!\mathit{ic}}
\def\Picardfunctor{\mathrm{Pic}}
\def\Deformationcategory{\mathcal{D}\!\mathit{ef}}

\setlength{\emergencystretch}{3em}
\begin{document}
\title[Stacks Project, Tag 0A2L]{277 --- The Brauer class of a degree-d Azumaya algebra is annihilated by d}
\maketitle
\noindent Copyright (C) 2005--2025 Johan de Jong. Stacks Project authors. Source: \href{https://stacks.math.columbia.edu/tag/0A2L}{Tag 0A2L}.

\noindent Editorial difficulty note (not author text). Difficulty H2: The proof must construct a canonical antisymmetrization section in the d-fold tensor power that glues despite line-bundle ambiguity in local splitting modules. Its right ideal gives a globally defined splitting module of the precise rank, yielding an endomorphism-algebra description without assuming d! is invertible..

\noindent Editorial provenance note (not author text). History: excerpt from revision \texttt{a0444\allowbreak 6e57e\allowbreak c1fbc\allowbreak 252a8\allowbreak 71afc\allowbreak ec775\allowbreak 2fb28\allowbreak 07b14}, etale-cohomology.tex, lines 9518--9574. Reference presentation was adapted; original-author context and core support blocks were retained or added. The mathematical wording is unchanged. Licensed under GNU FDL 1.2 or later; no invariant sections or cover texts. See ../licenses/COPYING. Context blocks reproduce author definitions and supporting results; they are not additional counted problems.

\begin{lemma}
\label{lemma-annihilated-by-degree}
\begin{reference}
Argument taken from \href{https://stacks.math.columbia.edu/bibliography}{Bibliography: Saltman-torsion}.
\end{reference}
Let $S$ be a scheme. Let $\mathcal{A}$ be an Azumaya algebra which is
locally free of rank $d^2$ over $S$. Then the class
of $\mathcal{A}$ in the Brauer group of $S$ is annihilated by $d$.
\end{lemma}

\begin{proof}
Choose an \'etale covering $\{U_i \to S\}$ and choose isomorphisms
$\mathcal{A}|_{U_i} \to \SheafHom(\mathcal{F}_i, \mathcal{F}_i)$
for some locally free $\mathcal{O}_{U_i}$-modules $\mathcal{F}_i$
of rank $d$. (We may assume $\mathcal{F}_i$ is free.) Consider the
composition
$$
p_i : \mathcal{F}_i^{\otimes d} \to
\wedge^d(\mathcal{F}_i) \to \mathcal{F}_i^{\otimes d}
$$
The first arrow is the usual projection and the second arrow is
the isomorphism of the top exterior power of $\mathcal{F}_i$ with
the submodule of sections of $\mathcal{F}_i^{\otimes d}$ which transform
according to the sign character under the action of the symmetric group
on $d$ letters. Then $p_i^2 = d! p_i$ and the rank of $p_i$ is $1$.
Using the given isomorphism
$\mathcal{A}|_{U_i} \to \SheafHom(\mathcal{F}_i, \mathcal{F}_i)$
and the canonical isomorphism
$$
\SheafHom(\mathcal{F}_i, \mathcal{F}_i)^{\otimes d} =
\SheafHom(\mathcal{F}_i^{\otimes d}, \mathcal{F}_i^{\otimes d})
$$
we may think of $p_i$ as a section of $\mathcal{A}^{\otimes d}$
over $U_i$. We claim that $p_i|_{U_i \times_S U_j} = p_j|_{U_i \times_S U_j}$
as sections of $\mathcal{A}^{\otimes d}$. Namely, applying
Lemma \href{https://stacks.math.columbia.edu/tag/0A2K}{lemma end unique up to invertible (Tag 0A2K)}
we obtain an invertible sheaf $\mathcal{L}_{ij}$ and a canonical isomorphism
$$
\mathcal{F}_i|_{U_i \times_S U_j} \otimes \mathcal{L}_{ij}
\longrightarrow
\mathcal{F}_j|_{U_i \times_S U_j}.
$$
Using this isomorphism we see that $p_i$ maps to $p_j$.
Since $\mathcal{A}^{\otimes d}$ is a sheaf on $S_\etale$
(Proposition \href{https://stacks.math.columbia.edu/tag/03OG}{proposition quasi coherent sheaf fpqc (Tag 03OG)}) we find a canonical
global section $p \in \Gamma(S, \mathcal{A}^{\otimes d})$. A local calculation
shows that
$$
\mathcal{H} =
\Im(\mathcal{A}^{\otimes d} \to \mathcal{A}^{\otimes d}, f \mapsto fp)
$$
is a locally free module of rank $d^d$ and that (left) multiplication
by $\mathcal{A}^{\otimes d}$ induces an isomorphism
$\mathcal{A}^{\otimes d} \to \SheafHom(\mathcal{H}, \mathcal{H})$.
In other words, $\mathcal{A}^{\otimes d}$ is the trivial element
of the Brauer group of $S$ as desired.
\end{proof}
\end{document}
